\section*{第 \ref{ch:b2:reproductive-hormones} 章——激素对生殖的调控}

\begin{solution}{exo:b2:reproductive-hormones:1}
下丘脑（GnRH，脉冲）$\to$ 垂体前叶（LH，FSH）$\to$
\omterm{def:b2:mammal-reproduction:testis}{睾丸}：LH 作用于\omterm{def:b2:mammal-reproduction:testis}{莱迪希细胞}（睾酮），FSH 作用于\omterm{def:b2:mammal-reproduction:testis}{支持细胞}
（\omterm{def:b2:mammal-reproduction:testis}{精子发生}，抑制素）。反馈：睾酮抑制 GnRH 和
LH；抑制素抑制 FSH。
\end{solution}

\begin{solution}{exo:b2:reproductive-hormones:2}
卵巢：\omterm{def:b2:reproductive-hormones:cycle}{卵泡期}，第 1--14 天，雌二醇来自生长中的
\omterm{def:b2:mammal-reproduction:ovary}{卵泡}；\omterm{def:b2:mammal-reproduction:ovary}{排卵}，第 14 天，LH 峰；\omterm{def:b2:reproductive-hormones:cycle}{黄体期}，第 15--28 天，
孕酮来自\omterm{def:b2:mammal-reproduction:ovary}{黄体}。子宫：月经，第 1--5 天，
（类固醇撤退）；增殖期，第 6--14 天（雌二醇）；
分泌期，第 15--28 天（孕酮）。
\end{solution}

\begin{solution}{exo:b2:reproductive-hormones:3}
每日服用的雌激素和孕激素通过负反馈使 LH 和 FSH 维持在低水平，
因此没有\omterm{def:b2:mammal-reproduction:ovary}{卵泡}被募集，雌二醇不会升高，
LH 峰也不会出现；孕激素还会使宫颈黏液变稠。
停药周内类固醇下降，在其作用下形成的\omterm{def:b2:reproductive-hormones:cycle}{子宫内膜}
脱落——这是撤药性出血，而非\omterm{def:b2:mammal-reproduction:ovary}{排卵}后的
月经。
\end{solution}

\begin{solution}{exo:b2:reproductive-hormones:4}
LH 和 FSH 升高数倍（睾酮和抑制素的反馈缺失）；
肌肉、性欲和附属性腺退化。补充睾酮可以
恢复第二性状并使 LH 回落，但不能恢复生育能力：
\omterm{def:b2:mammal-reproduction:testis}{睾丸}已经切除，即使是有\omterm{def:b2:mammal-reproduction:testis}{睾丸}的男性，外源
睾酮也会抑制 LH，进而使\omterm{def:b2:mammal-reproduction:testis}{精子发生}所需的
\omterm{def:b2:mammal-reproduction:testis}{睾丸}内睾酮减少。
\end{solution}

\begin{solution}{exo:b2:reproductive-hormones:5}
每天 $24/1.5 = 16$ 次脉冲；$24/4 = 6$。\omterm{def:b2:mammal-reproduction:ovary}{黄体}死亡时
脉冲仍然较慢，有利于 FSH：FSH 首先升高——
这恰是正确的激素，因为正是 FSH 募集下一批
\omterm{def:b2:mammal-reproduction:ovary}{卵泡}。
\end{solution}

\begin{solution}{exo:b2:reproductive-hormones:6}
\qty{200}{pg/mL} $= \qty{200}{ng/L}$: $200\times 10^{-9}/272 =
\qty{7.4e-10}{mol/L} = \qty{740}{pmol/L}$. \qty{15}{ng/mL} $=
\qty{15}{\micro g/L}$：$15\times 10^{-6}/314 = \qty{48}{nmol/L}$。在
一微升中：$7.4\times 10^{-16}$ 摩尔 $= 4.4\times 10^{8}$
个分子。
\end{solution}

\begin{solution}{exo:b2:reproductive-hormones:7}
\omterm{def:b2:mammal-reproduction:ovary}{排卵}在第 $35 - 14 = 21$ 天；易受孕期为第 18 至 22 天。
周期不规律时，\omterm{def:b2:reproductive-hormones:cycle}{卵泡期}及\omterm{def:b2:mammal-reproduction:ovary}{排卵}日
无法由以往周期推算，因而易受孕期
可能被漏算或错算。
\end{solution}

\begin{solution}{exo:b2:reproductive-hormones:8}
从第 21 天到第 26 天共 2.5 次倍增：$2^{2.5} = 5.7$ IU/L——足够。
若于第 25 天\omterm{def:b2:mammal-reproduction:implantation}{着床}，第 26 天为 $1.4$ IU/L；\omterm{def:b2:mammal-reproduction:ovary}{黄体}
已在退化，无法挽救，孕酮下降，
胚胎随月经丢失。\omterm{def:b2:mammal-reproduction:implantation}{着床}过晚是
早期流产的常见原因。
\end{solution}

\begin{solution}{exo:b2:reproductive-hormones:9}
垂体对 GnRH 的模式而非其量做出反应：
每小时一次的脉冲维持 LH 和 FSH，而同样的日剂量持续输注
则会完全抑制它们。这排除了简单的剂量–反应关系。
一个月的长效激动剂：经数日的激发效应后，LH 降至
接近零（\omterm{prop:b2:reproductive-hormones:pulses}{受体脱敏}），雌二醇降至\omterm{prop:b2:reproductive-hormones:overrides}{绝经}后
水平，\omterm{def:b2:reproductive-hormones:cycle}{子宫内膜}萎缩，不再出现\omterm{def:b2:reproductive-hormones:cycle}{月经周期}——
一种可逆的医学性\omterm{prop:b2:reproductive-hormones:overrides}{绝经}。
\end{solution}

\begin{solution}{exo:b2:reproductive-hormones:10}
无激素时：$R^* = 100/0.5 = 200$。持续给药：$100/6.5 = 15$。
一次两分钟的脉冲去除 $6\times 200\times(2/60) = 40$ 个受体
（\qty{20}{\%}）；随后一小时内缺口缩小为
原来的 $e^{-0.5}$ 倍，即其中 \qty{39}{\%} 得以恢复。受体池
稳定在每次脉冲的损失与每小时恢复量相等处：
每次脉冲前约 150 个受体，为最大值的
四分之三——约为持续给药水平的十倍。
\end{solution}

\begin{solution}{exo:b2:reproductive-hormones:11}
\omterm{prop:b2:reproductive-hormones:overrides}{绝经}：无\omterm{def:b2:mammal-reproduction:ovary}{卵泡}，故无雌二醇或抑制素，因而无
反馈：LH 和 FSH 居高不下。\omterm{prop:b2:mammal-reproduction:lactation}{催乳素}瘤：\omterm{prop:b2:mammal-reproduction:lactation}{催乳素}减慢了
GnRH 脉冲发生器，故 LH 偏低，\omterm{def:b2:mammal-reproduction:ovary}{卵泡}得不到驱动，
雌二醇偏低，无\omterm{def:b2:reproductive-hormones:cycle}{月经周期}。颗粒细胞瘤：
持续的高水平雌二醇通过负反馈抑制 LH（且
无成熟\omterm{def:b2:mammal-reproduction:ovary}{卵泡}，则无卵可排），因而
周期在雌二醇居高时停止——\omterm{def:b2:reproductive-hormones:cycle}{子宫内膜}也随之过度增生。
\end{solution}

\begin{solution}{exo:b2:reproductive-hormones:12}
雌二醇水平低或作用时间短时，会减少 GnRH 和 LH；高水平、持续的
雌二醇（高于 \qty{200}{pg/mL} 且持续超过 36 小时）会改变
kisspeptin--GnRH 系统和垂体的反应，
而垂体此时已储备了 LH，因此同一激素现在
引发一次 LH 爆发。\omterm{def:b2:mammal-reproduction:ovary}{排卵}必须全或无且定时：卵子
必须只释放一次，且只在某一天、\omterm{def:b2:mammal-reproduction:ovary}{卵泡}成熟时释放。若
采用旋钮式调节（LH 与雌二醇成比例），则将产生渐进、部分的
\omterm{def:b2:mammal-reproduction:ovary}{黄体}化，\omterm{def:b2:mammal-reproduction:ovary}{卵泡}可能半成熟时\omterm{def:b2:mammal-reproduction:ovary}{排卵}，或数个同时\omterm{def:b2:mammal-reproduction:ovary}{排卵}，
且不存在确定的易受孕日。
\end{solution}

\begin{solution}{pb:b2:reproductive-hormones:1}
\textbf{1.} \omterm{def:b2:reproductive-hormones:cycle}{卵泡期}，第 1--14 天：雌二醇上升，孕酮
较低。\omterm{def:b2:reproductive-hormones:cycle}{黄体期}，第 15--28 天：孕酮较高。
\textbf{2.} LH 峰于第 14 天达到峰值（第 13 天开始）；
\omterm{def:b2:mammal-reproduction:ovary}{排卵}约在开始后 36 小时，即第 15 天。
\textbf{3.} 第 15 至 28 天：14 天，即\omterm{def:b2:mammal-reproduction:ovary}{黄体}的
寿命。
\textbf{4.} 孕酮使体温调定点升高
\qty{0.3}{\celsius}至\qty{0.4}{\celsius}；此升高在\omterm{def:b2:mammal-reproduction:ovary}{排卵}后一两天出现，
因而只能证实\omterm{def:b2:mammal-reproduction:ovary}{排卵}已经发生，却
无法预先提示\omterm{def:b2:mammal-reproduction:ovary}{排卵}。
\textbf{5.} 第 11 至 14 天（210、260、310、220）：持续三至四天，
超过开关所需的 36 小时。
\textbf{6.} 第 12 至 16 天。
\textbf{7.} $310\times 10^{-9}/272 = \qty{1.14e-9}{mol/L} =
\qty{1140}{pmol/L}$; $16\times 10^{-6}/314 = \qty{51}{nmol/L}$.
\textbf{8.} $65/8 = 8$。半衰期为一小时，LH 在分泌停止后
每小时减半，因此一天后已接近基线水平。
\textbf{9.} \omterm{def:b2:mammal-reproduction:ovary}{排卵}约在第 22 天，周期约 36 天。
\textbf{10.} 第 22 天\omterm{def:b2:mammal-reproduction:implantation}{着床}；至第 27 天（五天，2.5
次倍增）约为 \qty{5.7}{IU/L}。
\textbf{11.} 第 27 天孕酮降至 \qty{2}{ng/mL}，说明
\omterm{def:b2:mammal-reproduction:ovary}{黄体}按计划退化，未受任何挽救。
\textbf{12.} 孕酮约为 \qty{20}{ng/mL} 并且上升，
雌二醇约为 \qty{200}{pg/mL}。
\textbf{13.} 周期末，雌二醇、孕酮和
抑制素一同下降，FSH 摆脱了其反馈抑制；随着\omterm{def:b2:mammal-reproduction:ovary}{卵泡}群
生长，雌二醇和抑制素上升，FSH 下降，只有
FSH 受体最多（颗粒细胞最多）的\omterm{def:b2:mammal-reproduction:ovary}{卵泡}
仍能在下降的 FSH 作用下继续生长，其余\omterm{def:b2:mammal-reproduction:ovary}{卵泡}死亡。
\textbf{14.} 无激素时 $R^* = 200$；在
持续激素作用下 $100/6.5 = 15$。
\textbf{15.} $6\times 200\times 2/60 = 40$ 个受体，\qty{20}{\%}；
时间常数 $1/d = \qty{2}{h}$，故缺口的 \qty{39}{\%} 在
一小时内恢复；受体池稳定在约 150 个受体，为最大值的
四分之三。
\textbf{16.} 脉冲使垂体维持在约 150 个受体，它
对每次脉冲以 LH 应答；持续输注使受体池降至
15，同样的细胞便不再应答，无论剂量多少；脉冲则恢复
受体池及应答。
\textbf{17.} 第一天：LH 和睾酮升高（激动剂
刺激仍存在的受体——即激发效应）。一个月后：
\omterm{prop:b2:reproductive-hormones:pulses}{受体脱敏}，LH 接近零，睾酮降至去势水平，
这正是治疗目标。
\textbf{18.} FSH 升高（无抑制素），LH 保持正常；更多\omterm{def:b2:mammal-reproduction:ovary}{卵泡}
每个周期都能通过选择存活，出现多次\omterm{def:b2:mammal-reproduction:ovary}{排卵}和
异卵双胞胎。
\textbf{19.} 在自然周期中，FSH 下降使除优势\omterm{def:b2:mammal-reproduction:ovary}{卵泡}外的所有\omterm{def:b2:mammal-reproduction:ovary}{卵泡}
死亡；持续供应的 FSH 使整个\omterm{def:b2:mammal-reproduction:ovary}{卵泡}群
生长至成熟。
\textbf{20.} 十个\omterm{def:b2:mammal-reproduction:ovary}{卵泡}在早期产生的雌二醇远超 \qty{200}{pg/mL}，
早期即可引发过早的 LH 峰，
并在取卵前\omterm{def:b2:mammal-reproduction:ovary}{排卵}；拮抗剂阻断 GnRH，故
不会出现 LH 峰。
\textbf{21.} 表中 LH 峰于第 14 天达到峰值，\omterm{def:b2:mammal-reproduction:ovary}{排卵}
在峰开始约 36 小时后出现；hCG 替代了
LH 峰，35 小时取卵时卵母细胞已成熟
但仍在\omterm{def:b2:mammal-reproduction:ovary}{卵泡}内。
\textbf{22.} $10\times 0.7 = 7$ 个胚胎，$7\times 0.4 = 2.8$
个\omterm{def:b2:mammal-reproduction:implantation}{囊胚}；$1 - 0.65^{2} = 0.58$。
\textbf{23.} 持续的雌激素和孕激素使 FSH 和 LH 维持低水平：无
\omterm{def:b2:mammal-reproduction:ovary}{卵泡}生长，雌二醇从不升至阈值，无 LH 峰，
无\omterm{def:b2:mammal-reproduction:ovary}{黄体}，孕酮亦不升高。
\textbf{24.} 须给予睾酮以维持第二性状
和性欲；但它无法恢复\omterm{def:b2:mammal-reproduction:testis}{莱迪希细胞}所提供的
高\omterm{def:b2:mammal-reproduction:testis}{睾丸}内浓度，因此\omterm{def:b2:mammal-reproduction:testis}{精子发生}仍被
抑制——这正是目的。
\textbf{25.} \omterm{def:b2:mammal-reproduction:ovary}{排卵}在第 15 天；\omterm{def:b2:reproductive-hormones:cycle}{黄体期} 14 天；易受孕期
为第 12--16 天；LH 峰由雌二醇触发，条件为
雌二醇高于 \qty{200}{pg/mL} 且持续超过 36 小时，于第 11--14 天满足。
\end{solution}
